% Source: https://artofproblemsolving.com/community/c3h1486063_ahsme_1957_problem
% Stable identity: aops-topic-1486063-q01
Draw $MB$. Since $AB$ is a diameter, $\angle AMB=90^\circ$; and $\angle AOP=90^\circ$ because the diameters are perpendicular.

\begin{asy}
usepackage("amsmath");
size(5cm);
defaultpen(fontsize(10pt));
path ra(pair P,pair u,pair v,real s=0.2){return P+s*unit(u)--P+s*unit(u)+s*unit(v)--P+s*unit(v);}
pair O=(0,0),A=(-1,0),B=(1,0),C=(0,1),D=(0,-1);
pair M=dir(-35); pair P=extension(A,M,C,D);
draw(circle(O,1),gray);
draw(A--B,linewidth(0.8)); draw(C--D,linewidth(0.8));
draw(A--M,blue+linewidth(0.9)); draw(M--B,red+dashed);
draw(ra(O,B-O,C-O,0.09)); draw(ra(M,A-M,B-M,0.09));
dot(O);dot(A);dot(B);dot(C);dot(D);dot(M);dot(P);
label("$A$",A,W); label("$B$",B,E); label("$C$",C,N); label("$D$",D,S); label("$M$",M,SE);
label("$O$",O,NW); label("$P$",P,NE);
label("\mbox{}",(min(currentpicture,user=true)+max(currentpicture,user=true))/2,black); // keeps the text after the figure black when LaTeX typesets the labels inline
\end{asy}

Triangles $AOP$ and $AMB$ are both right triangles and share the angle at $A$, so they are similar:
\[
\dfrac{AP}{AB}=\dfrac{AO}{AM}\quad\Longrightarrow\quad AP\cdot AM=AO\cdot AB .
\]
The answer is $\boxed{\textbf{(B)}}$. (For example, if $M=B$, then $P=O$ and both sides equal $AO\cdot AB$.)
